% !TEX program = lualatex
\documentclass[11pt,a4paper,luatex]{jlreq}

\usepackage{luatexja-fontspec}
\usepackage{geometry}
\geometry{margin=25mm}

\usepackage{amsmath,amssymb,amsthm}
\usepackage{mathtools}
\usepackage{tikz-cd}
\usepackage{enumitem}
\usepackage{hyperref}

\theoremstyle{definition}
\newtheorem{definition}{定義}[section]

\theoremstyle{plain}
\newtheorem{proposition}[definition]{命題}
\newtheorem{theorem}[definition]{定理}

\renewcommand{\proofname}{論証}
\renewcommand{\qedsymbol}{$\blacksquare$}

\title{\textbf{関係基底の存在論と意味の場の圏論的定式化：} \\
\large 言語的多義性による構造崩壊とユニバースの内包としての現象}

\author{太田 寛（Hiroshi Ota）\thanks{10.5281/zenodo.24267213: \url{https://doi.org} (All Versions: \url{https://doi.org}) \\ Homepage: \url{https://kae-made.jp}, Research Notes: \url{https://note.com}} \\
\small Knowledge \& Experience / 独立研究者 \\
\small \texttt{master@kae-made.jp}}

\date{\today}

\begin{document}

\maketitle

\begin{abstract}
従来の言語哲学および一般的なモデル化技法の多くは、対象（モノ）を基底とし、関係をその間に付随する二次的な属性とみなしてきた。しかし、この対象基底の枠組みでは、日常言語が持つ「多義性」によって意味ネットワークが複雑に絡み合いもつれ合って、構造が自明化する問題を本質的に説明できない。本稿では、存在（モノ）と意味的リンクの圏から出発し、その双対的な視点である「関係（リンク）を実在の主題とするスパン構造」こそが意味成立の本質であることを示す。さらに、単語によるラベリング関手が持つ多義性によって異なる意味のリンクが縮退し、恒等射と同型になることで意味が崩壊するメカニズムを定式化する。この崩壊を回避するため、人間は予め規定された概念の物差し（スキーマ圏）によって世界を分割し、スキーマに従うインスタンス圏（意味の場）と、場と場の対応を規定するブリッジからなる「ユニバース」を構成する。最終的に、私たちが「現象」や「観測」と呼ぶ認識は、このユニバースの「内包」として現れるものであり、観測された内包から元のユニバース全体の構造を一意に逆算することは原理的に不可能であることを論証する。
\end{abstract}

\section{導入}

私たちが世界を認識し、思考し、言語によって記述するとき、そこには常に「実在する世界」と「認識された現象」の間の構造的なギャップが存在する。カントが『純粋理性批判』\cite{kant1781} において、人間が認識できるのは「物自体」ではなく認識の枠組みを通して立ち現われる「現象」に限られる、と論じたように、実在の全体構造と私たちが把握する現象の間には、本質的な非対称性が存在する。また、フレーゲの文脈原理 \cite{frege1892} や前期ウィトゲンシュタインの写像理論 \cite{wittgenstein1921} に代表される従来の分析哲学・存在論は、「世界は対象（モノ）の総体であり、意味とは対象同士の配置や関係によって生じる」という対象基底の視座に立ってきた。

しかし、この視座には根本的な困難がある。私たちが現象を記述するために用いる単語（ラベル）は本質的に「多義性」を帯びているため、単一のネットワーク上でモノとモノを言葉で結びつけると、異なる文脈のリンクが複雑に絡み合い、ついには本来異なるはずの関係同士が区別不能（自明）となって意味が崩壊してしまうのである。これは、マルクス・ガブリエルが新実在論において「すべてを包み込む単一の『世界』は存在せず、実在は無数の異なる『意味の場（Fields of Sense）』に現れる」と指摘した事態の、構造的・数理的な帰結にほかならない \cite{gabriel2013,gabriel2015}。

本稿の目的は、日常的な言語理解から物理学・自然科学における現象の理論的記述に至るまで、人間が世界を分節化して捉える「一般的なモデル化技法」の共通基盤を、圏論 \cite{maclane1998} によって定式化し、それ単体で完結する以下の3つの命題を論証することである。

\begin{enumerate}
    \item \textbf{区別と関係基底への双対的転換：} 認識における「区別」とは対象に固有の恒等射（自己同一性）を認めることであり、指差し可能な実在は「モノ」「特徴値」「意味的リンク」に分類される。このうち、意味を成す実在の本質は「意味的リンク（関係）」そのものであり、リンクが自己同一性に潰れず差異を保つためには、リンクを対象とした、左右の役割へ射影する双対的スパン構造が要請される。
    \item \textbf{エポケーとスキーマ分割によるユニバースの構成：} 言語の多義性による意味の崩壊を防ぐため、人間は無意識に当てている単一の物差しを一旦停止（エポケー）し、複数のスキーマ圏（概念の物差し）によって世界を分割する。これにより、それぞれのインスタンス圏（意味の場）とそれらを結ぶブリッジからなる総体（ユニバース）が定式化される。
    \item \textbf{現象の「内包」性と逆算不可能性：} 私たちが観測と呼ぶ現象は、このユニバースの「内包」として切り出された認識の束であり、観測結果から元のユニバースの全体構造を一意に逆算することは原理的に不可能である。
\end{enumerate}

\section{認識における「区別」と存在の圏論的定式化}

\subsection{「区別」と恒等射、および指差し可能な3対象}

私たちが現象の中から何かを「区別する（個体化する）」とは、他との境界を引き、それに固有の自己同一性を認める行為にほかならない（クワインの「同一性なくして実体なし」の原則 \cite{quine1969}）。圏論の枠組みにおいて、この「区別され、指差し可能な存在」は、それ自身を指し示す\textbf{恒等射}（identity morphism）$\mathrm{id}$ を持つ対象として厳密に定式化される（圏論において対象を恒等射そのものと同一視する「射のみの定義」が成り立つように、固有の恒等射の確立こそが実在の根拠となる \cite{maclane1998}）。

現象の中で指差し可能な実在（恒等射を持つ対象）は、その構造的役割によって次の3つに厳密に分類される。

\begin{definition}[指差し可能な実在の3分類]
現象の中で指差し可能な実在は、次の3つに分類される。
\begin{enumerate}
    \item \textbf{モノ（実体 $A, B, \dots$）：} 現象の中で独立した個体として立ち現れ、特徴や関係の結節点となる対象。固有の恒等射 $\mathrm{id}_{A}$ を持つ。
    \item \textbf{特徴値（属性値 $V_1, V_2, \dots$）：} モノに付随する性質や状態（例：「赤い」「温度が高い」など）。特徴値 $V$ も指差し可能な対象として恒等射 $\mathrm{id}_{V}$ を持つが、その構造はモノ $A$ に対して一方向の従属的な射 $\alpha: A \to V$ で結ばれるのみであり、それ自体が左右に別のモノを引き連れて関係を拡張する力を持たない。
    \item \textbf{意味的リンク（関係 $R$）：} モノとモノの間に成立する意味的なつながりそのもの。
\end{enumerate}
\end{definition}

まず、素朴な認識として、モノ $A, B$ を対象とし、モノとモノの間に張られる意味的リンクを射とする圏 $\mathcal{C}_{\mathrm{inst}}$ を考えることができる。ここで意味的リンクは、一方向の射で完結する特徴値とは異なり、常に往路 $f: A \to B$ と復路 $g: B \to A$ の\textbf{双方向ペア}として現れる。なぜなら、あるモノ $A$ が別のモノ $B$ に対して意味を持つとき、そこには必ず $B$ から $A$ を参照する逆向きの文脈が随伴的に規定されなければ、関係としての輪郭が閉じないからである。

\subsection{双対的な視点の転換：リンク中心のスパン構造}

しかし、モノを対象（主役）とし、リンクを単なるモノ間の矢印（従属物）とみなす $\mathcal{C}_{\mathrm{inst}}$ の見方では、意味が成立する本質的な条件を見落とすことになる。一方向の矢印で閉じる「特徴値」とは決定的に異なり、真に意味のネットワークを織り成す実在の核は「意味的リンク（関係）$R$ そのもの」だからである。

リンク $R$ を第一級の実在（対象）として捉え直すと、$\mathcal{C}_{\mathrm{inst}}$ とは対象と射の役割が反転した\textbf{双対的な構造}が立ち上がる。

\begin{definition}[リンク中心の表現とスパン構造]
意味的リンク $R$ そのものを対象とみなすとき、$R$ は実在である以上、自身を指す恒等射 $\mathrm{id}_R: R \to R$ を持つ。同時に、$R$ からは関係の両端に位置するモノ $A, B$ への射影 $p_1: R \to A$ および $p_2: R \to B$ が伸びる。数学的形式として、これは $R$ を頂点とし左右にモノを引き連れる\textbf{スパン（Span）} \cite{benabou1967}
\begin{equation*}
A \xleftarrow{p_1} R \xrightarrow{p_2} B
\end{equation*}
である。同時に、論理命題の観点からは、この2つの射影が左右の役割（フレーズ）を担うことで「$A$ と $B$ が関係 $R$ にある」という一つの命題文（述語構造）を成立させる基盤となっている。
\end{definition}

$\mathcal{C}_{\mathrm{inst}}$ において「モノ $A$ と $B$ の間に双方向リンクがある」と観測していた現象は、この本質的なスパン構造を端点であるモノの側から眺めた射影に過ぎない。

\begin{proposition}[意味成立の非自明性条件]
意味的リンク $R$ が意味を成すとは、$R$ が自身の恒等射 $\mathrm{id}_{R}$ と同型な自明な点に縮退せず、左右への射影 $p_1, p_2$ によって異なる役割（差異）を保持していることである。
\end{proposition}

もしリンクが左右への非自明な射影のペアを失い、単一の閉じた点と区別がつかなくなれば、$R$ は単なる $\mathrm{id}_{R}$ へと還元され、関係としての意味を喪失する。

\begin{figure}[htbp]
\centering
\begin{minipage}{0.42\textwidth}
\centering
\textbf{対象基底（$\mathcal{C}_{\mathrm{inst}}$）}\\[1.2em]
\begin{tikzcd}[column sep=large]
A \arrow[r, shift left=1ex, "f"] & B \arrow[l, shift left=1ex, "g"]
\end{tikzcd}
\end{minipage}
\hfill
\begin{minipage}{0.52\textwidth}
\centering
\textbf{関係基底（スパン構造）}\\[1.2em]
\begin{tikzcd}[column sep=large]
A & R \arrow[l, swap, "p_1"] \arrow[r, "p_2"] & B
\end{tikzcd}
\end{minipage}
\caption{対象基底の双方向射ペア（$A \rightleftarrows B$）から、意味的リンク $R$ を主題（対象）として左右の役割へ射影するリンク中心のスパン構造（$A \xleftarrow{p_1} R \xrightarrow{p_2} B$）への双対的転換。}
\label{fig:span_shift}
\end{figure}

\section{言語的ラベリングの多義性による意味の崩壊}

通常の言語活動や思考において、個別の実在にその都度固有の記号を割り振ることはせず、既存の「単語（ラベル）」を用いてモノや関係を指示する。このプロセスは、実在の圏から言語ラベルの圏 $\mathcal{W}$ への射影（ラベリング関手）として記述できる。

\begin{definition}[ラベリング関手と多義性]
実在のネットワークを成す圏 $\mathcal{E}$ から、単語ラベルを対象とする圏 $\mathcal{W}$ への関手 $L: \mathcal{E} \to \mathcal{W}$ を考える。ある単語 $w \in \mathcal{W}$ が\textbf{多義性}を持つとは、$L$ による逆像（ファイバー）$L^{-1}(w)$ の中に、本来異なる文脈・異なるスパン構造に属する複数の実在対象が混在している状態を指す。
\end{definition}

たとえば、「橋（bridge）」という単一のラベルは、物理的構造物としての橋（川と岸をつなぐ関係）と、比喩的な仲介（人と人をつなぐ関係）の双方に適用される。このように多義的なラベリング関手 $L$ を通して世界を単一のネットワーク上に記述しようとすると、ネットワークは複雑な蜘蛛の巣状の絡み合いを呈する。

\begin{theorem}[多義性による恒等射への縮退と意味の崩壊]
多義的なラベリング関手 $L$ によって異なる文脈の対象が同一のラベル $w$ に畳み込まれる（商圏化される）とき、本来異なる対象間を結んでいた意味的リンク $R$ の両端が同一視され、自己ループを形成する。さらに文脈の差異が失われることで、その射はラベル $w$ の恒等射 $\mathrm{id}_{w}$ と同型に縮退し、命題 2.3 の非自明性条件が破れて意味が崩壊する。
\end{theorem}

すなわち、言葉の多義性を無秩序に許容した単一のネットワークモデルでは、あらゆる意味的リンクが最終的には自明な恒等射へと潰れてしまい、現象を正しく記述することが不可能となる。

\begin{figure}[htbp]
\centering
\begin{tikzcd}[column sep=large, row sep=large]
\mbox{\textbf{実在（$\mathcal{E}$）：}} &
A_1 & R_1 \arrow[l, swap, "p_1"] \arrow[r, "p_2"] \arrow[dd, dashed, "L"] & B_1 \quad\;\;
A_2 & R_2 \arrow[l, swap, "q_1"] \arrow[r, "q_2"] \arrow[ddll, dashed, swap, "L"] & B_2 \\
\\
\mbox{\textbf{ラベル（$\mathcal{W}$）：}} &
& w \arrow[loop right, "\mathrm{id}_w\mbox{（縮退）}"] & & &
\end{tikzcd}
\caption{言語的ラベリングの多義性による意味の崩壊（定理 3.2）。実在の圏 $\mathcal{E}$ における異なるスパン構造 $R_1, R_2$ が、多義的なラベル $w \in \mathcal{W}$ のファイバー $L^{-1}(w)$ へと畳み込まれることで、恒等射 $\mathrm{id}_w$ と同型な自明な自己ループへと縮退する。}
\label{fig:polysemy_collapse}
\end{figure}


\section{意味の場とブリッジによるユニバースの構成、および現象学的エポケー}

\subsection{存在の形式としてのスキーマ・意味の場・ブリッジ、および現象学的エポケー}

第3章で示した通り、異なる文脈に属する関係（意味的リンク）を単一の同質な空間に押し込めようとすると、リンク同士が混線・縮退して恒等射と区別がつかなくなり、構造そのものが崩壊する。したがって、世界が非自明な差異を保ったまま実在するためには、観測者の有無にかかわらず、実在の側があらかじめ文脈ごとに分節化された構造をとっていなければならない。

この分節化された存在の形式を規定するのが、構造の型である「スキーマ圏」と、それに従って具体的な実在が配置された「インスタンス圏（意味の場）」である。これを人間の認識の側に当てはめてみれば、私たちが現実の思考において意味の崩壊に陥らないのは、現象に対して無秩序に単語を貼り付けるのではなく、「予め分類された概念という物差し（スキーマ）」を当てはめて世界を分節化しているからにほかならない。

\begin{definition}[スキーマ圏と意味の場]
特定の文脈や法則体系における概念クラス（モノの型）、属性型（特徴値の型）、およびリレーション型（関係の型）によって規定される圏を\textbf{スキーマ圏} $\mathcal{S}_{i}$ と定義する。そして、スキーマ $\mathcal{S}_{i}$ の構造に従って具体的なモノ・特徴値・意味的リンクが配置された実在の局所的な状態を、関手 $I_i: \mathcal{S}_i \to \mathbf{Set}$（あるいはその圏論的表現としての\textbf{インスタンス圏} $\mathcal{I}_{i}$）として定式化する \cite{lawvere1963,spivak2012}。
\end{definition}

ここで重要なのは、スキーマ $\mathcal{S}_{i}$ とインスタンス $\mathcal{I}_{i}$ が「構文（関係を主題とする型の規定）」と「意味（モノを主題とする具体的な場）」という双対的な関係（関手意味論の双対性 \cite{lawvere1963}）にある点である。こうして各スキーマ $\mathcal{S}_{i}$ ごとに切り出された局所的なインスタンス圏 $\mathcal{I}_{i}$ こそが、マルクス・ガブリエルの言う\textbf{「意味の場（Field of Sense）」}の圏論的定式化である \cite{gabriel2015}。

しかし、日常の素朴な認識（フッサール現象学における「自然的態度」\cite{husserl1913}）においては、人間は特定の慣習的な物差し（単一のスキーマ）を無意識のうちに現象へ当てはめてしまい、それが唯一の実在であるかのように振る舞ってしまう。その結果、異なる意味の場が重なり合ってファイバーの衝突が起き、意味の崩壊が生じる。このとき観測者の側で機能するのが、現象学における\textbf{エポケー（Epoch\'e：判断中止）}に相当する操作である。エポケーとは、無意識に当てていた既存の単一の物差しの適用を一旦停止（括弧入れ）し、対象を単一ラベルの束として受け入れる態度を保留することを意味する。この一時停止によって、私たちを、現象が単一の枠組みへ押し込められてしまう事態から解放し、実在が本来備えている複数の分節構造（異なるスキーマ圏 $\{\mathcal{S}_i\}$ と意味の場 $\{\mathcal{I}_i\}$ の配置）へと、視点を同調させることが可能となる。各意味の場 $\mathcal{I}_{i}$ の内部では文脈の混線が排除され、すべての意味的リンクが恒等射と区別される健全な非自明性を保持することになる。

従来の現象学がエポケーを「超越論的主観性（純粋意識）」という単一の究極的基盤への還元として位置づけていたのは、多様な現象を唯一の中心へと統合しようとする試みの帰結にほかならない \cite{husserl1913}。しかし、異なる文脈を単一の空間に統合しようとすれば構造的破綻が生じるため、そこには常に経験を超えた形而上学的装置を仮定する必要が生じてしまう。これに対し本稿は、エポケーを単一の基盤への還元ではなく、無意識に固定していた単一の物差しを一旦停止し、実在が本来備えている複数の並立する意味の場 $\{\mathcal{I}_i\}$ とスキーマ圏 $\{\mathcal{S}_i\}$ の構造へと視点を解放する操作として再定義するものである。

\subsection{意味の場をつなぐブリッジとユニバース}

このように見れば、言語やモデル化において現れる多義性とは、単なる人間側の表記上の不備やノイズではない。実在そのものが複数の意味の場に分節化され、かつ異なる場に属する実体が対応関係をもって並存しているという存在の形式そのものから避けようもなく生じる事態なのである。したがって、分割された意味の場同士を再び関係づける客観的な構造が必要となる。

\begin{definition}[ブリッジとユニバース]
異なるスキーマ $\mathcal{S}_i, \mathcal{S}_j$ に基づく2つの意味の場 $\mathcal{I}_{i}$ と $\mathcal{I}_{j}$ の間にある対応関係を、媒介構造 $\mathcal{B}_{ij}$ を頂点とするスパン \cite{fong2019}
\begin{equation*}
\mathcal{I}_i \leftarrow \mathcal{B}_{ij} \rightarrow \mathcal{I}_j
\end{equation*}
として定義し、これを\textbf{ブリッジ（Bridge）}と呼ぶ。
\end{definition}

ここで極めて重要なのは、\textbf{ブリッジ $\mathcal{B}_{ij}$ が対応づける対象は各意味の場における「モノ」および「特徴値」に限定され、「意味的リンク」そのものを直接対応づけることはない}という点にある。意味的リンク $R$ は各スキーマ $\mathcal{S}_{i}$ 固有の文脈（スパン構造）に強く依存する関係性であり、異なる場のリンク同士を直接同一視すれば再び構造崩壊を招くためである。したがって、異なる意味の場同士は、あくまで端点となるモノや特徴値の対応（ブリッジ）を介してのみ接続され、意味的リンクの構造そのものは、各意味の場の内部に閉じている。

そして、すべてのスキーマ圏 $\{\mathcal{S}_i\}$、それに従う意味の場（インスタンス圏）$\{\mathcal{I}_i\}$、およびモノと特徴値のみを結ぶブリッジ $\{\mathcal{B}_{ij}\}$ が統合された全体構造を\textbf{ユニバース} $\mathcal{U}$ と定義する。このユニバース $\mathcal{U}$ こそが、観測者の存在に先立って成り立つ客観的な「存在の形式」である。

\begin{figure}[htbp]
\centering
\begin{tikzcd}[column sep=large, row sep=large]
& \mathcal{B}_{ij} \arrow[dl, swap, "\mbox{モノ・特徴値のみ}"] \arrow[dr, "\mbox{モノ・特徴値のみ}"] & \\
\mathcal{I}_i\mbox{（意味の場 } i\mbox{）} \arrow[rr, dashed, no head, "\mbox{意味的リンク } R \mbox{ は内部に閉じる}"'] & & \mathcal{I}_j\mbox{（意味の場 } j\mbox{）} \\
\mathcal{S}_i\mbox{（スキーマ } i\mbox{）} \arrow[u, "I_i"] & & \mathcal{S}_j\mbox{（スキーマ } j\mbox{）} \arrow[u, swap, "I_j"]
\end{tikzcd}
\caption{ユニバース $\mathcal{U}$ の全体構造（定義 4.2）。異なる意味の場 $\mathcal{I}_i$ と $\mathcal{I}_j$ は、それぞれ対応するスキーマ $\mathcal{S}_i, \mathcal{S}_j$ からのインスタンス関手 $I_i, I_j$ によって規定される。意味的リンク $R$ は各場の内部に厳密に閉じており、場と場の間の対応は「モノ」と「特徴値」のみを媒介するブリッジスパン $\mathcal{I}_i \leftarrow \mathcal{B}_{ij} \rightarrow \mathcal{I}_j$ によってのみ接続される。}
\label{fig:universe_bridge}
\end{figure}

\section{結論：ユニバースの内包としての現象と逆算不可能性}

以上の定式化により、実在と認識の関係は次のように結論づけられる。

\begin{proposition}[実在の形式と「内包」としての現象]
この世界の実在とは、スキーマを伴うインスタンスの意味の場の集合と、それらの間のブリッジからなるユニバース $\mathcal{U}$ の住人である。そして、私たちが普段「現象」や「観測」と呼んでいる認識は、このユニバース $\mathcal{U}$ の内部に折りたたまれ、特定のスキーマと文脈のフィルターを通して切り出されたユニバースの\textbf{「内包（Intension）」}にほかならない。
\end{proposition}

最後に、この構造から導かれる最も重要な帰結である「逆算不可能性」を証明する。

\begin{theorem}[観測からの逆算不可能性]
観測された現象（ユニバースの内包）のデータのみから、元のユニバース $\mathcal{U}$ 全体の構造（どのスキーマ群 $\mathcal{S}_{i}$ とどのブリッジ $\mathcal{B}_{ij}$ の組み合わせによってその現象が生成されたか）を一意に逆算することは不可能である。
\end{theorem}

\begin{proof}
本定理を論証するため、実在のユニバース $\mathcal{U}$ から、私たちが認識する「現象（内包）」が生成される観測プロセスを、圏論的な射影（関手）として定式化し、その逆写像が存在し得ないこと（単射性の破れ）を3つの段階に分けて示す。

\paragraph{第1段階：観測プロセスの「射影（忘却関手）」としての定式化．}
実在のユニバース $\mathcal{U}$ は、互いに文脈を異にする複数のスキーマ圏 $\{\mathcal{S}_i\}$、それらに従うインスタンス圏（意味の場）$\{\mathcal{I}_i\}$、および各場を結ぶブリッジ $\{\mathcal{B}_{ij}\}$ が統合された高次の構造体である。これに対し、人間が日常の思考や科学的手続きによって世界を「観測」し、モデルや言語として把握する行為は、実在のユニバース $\mathcal{U}$ に対して特定の主観的スキーマや語彙ラベルの圏 $\mathcal{W}$ という「フィルター」を適用し、観測可能な側面だけを切り出す操作に相当する。

圏論の枠組みにおいて、この観測プロセスは、ユニバース $\mathcal{U}$ の豊かな構造から一部の情報（文脈の差異や高次の関係構造）を削ぎ落として低次の表現へと写し取る\textbf{射影関手（忘却関手：Forgetful Functor）} $\Pi: \mathcal{U} \to \mathcal{O}$（ここで $\mathcal{O}$ は観測された現象の圏）として定式化される。命題 5.1 で述べた通り、私たちが経験する「現象」とは、この射影関手 $\Pi$ によって切り出されたユニバースの内包（Intension）にほかならない。

\paragraph{第2段階：射影過程における2つの構造的情報損失．}
観測の射影関手 $\Pi$ が原像（ユニバース $\mathcal{U}$）の構造を完全に保存する（忠実充満である）ならば、観測結果 $\Pi(\mathcal{U})$ から元の $\mathcal{U}$ を一意に復元することが可能となる。しかし、第3章および第4章で明らかにした構造的制約により、射影関手 $\Pi$ の適用過程では、次の2つの不可逆な情報損失（縮退）が必然的に発生する。
\begin{enumerate}
    \item \textbf{言語ラベリングの多義性によるファイバーの縮退：} 観測結果を記号や単語ラベルの圏 $\mathcal{W}$ へ落とし込む際、ラベリング関手 $L$ の非単射性（多義性）により、本来異なる意味の場 $\mathcal{I}_i, \mathcal{I}_j$ に属する別個の実在対象や意味的リンクが、同一のラベル $w$ の逆像（ファイバー）$L^{-1}(w)$ へと重ね合わされる。この畳み込み（商圏化）によって、異なるスパン構造が持っていた固有の差異は相殺され、観測面上では区別不能な単一の現れへと縮退する。
    \item \textbf{ブリッジを介した「意味的リンク」の非保存性：} 第4.2節で定義した通り、異なる意味の場 $\mathcal{I}_{i}$ と $\mathcal{I}_{j}$ を結ぶブリッジ $\mathcal{B}_{ij}$ は、端点となる「モノ」と「特徴値」の対応関係のみを媒介し、各スキーマ固有の文脈に依存する「意味的リンク $R$」そのものを直接対応づけることはない。したがって、ある特定の意味の場 $\mathcal{I}_{i}$ を窓としてユニバースを観測した場合、ブリッジの向こう側にある別の意味の場 $\mathcal{I}_{j}$ に属するモノの存在は間接的に察知できたとしても、その場においてそれらのモノがいかなる意味的リンク（関係構造）によって織り成されているかという情報は、$\Pi$ によって完全に脱落（忘却）する。
\end{enumerate}

\paragraph{第3段階：単射性の破れと逆写像の非存在．}
以上の2つの構造的情報損失により、射影関手 $\Pi: \mathcal{U} \to \mathcal{O}$ は単射（忠実）ではない。すなわち、内部に異なるスキーマ分割や異なる意味的リンクの構造を持つ2つの異なるユニバース部分構造 $\mathcal{U}_1 \neq \mathcal{U}_2$ を想定したとき、それぞれの関係構造の差異がラベリングの多義性によって縮退し、かつブリッジの向こう側のリンク構造が忘却されることで、観測される現象（内包）としては完全に同一の像を結ぶケースが必然的に存在する：
\begin{equation*}
\mathcal{U}_{1} \neq \mathcal{U}_{2} \quad \text{でありながら} \quad \Pi(\mathcal{U}_{1}) \cong \Pi(\mathcal{U}_{2}).
\end{equation*}
数学的および論理的に、異なる原像に対して同一の像を返す写像（非単射な関手）には、像から原像を一意に決定する逆写像 $\Pi^{-1}$ が存在しない。したがって、観測された現象（ユニバースの内包）のデータのみを出発点として、それを生み出した実在のユニバース $\mathcal{U}$ 全体の構造を一意に逆算・再構成することは原理的に不可能である。
\end{proof}

\subsection*{まとめ}

本稿では、カントの認識論的枠組みを起点としつつ、日常の思考から自然科学の理論記述に至る「一般的なモデル化技法」の根底にある構造を圏論的に定式化した。認識における「区別」を恒等射の確立として定義し、指差し可能な実在から「特徴値（一方向の射）」と「意味的リンク（左右にモノを引き連れるスパン構造）」を厳密に区別した上で、関係を実在の本質とする双対的視点から言語の多義性がもたらす意味崩壊のメカニズムを明らかにした。さらに、人間が現象学的エポケーによって無意識の物差しを一旦停止し、複数のスキーマを用いて世界を「意味の場（Fields of Sense）」へと分割し、モノと特徴値のみを結ぶブリッジによってそれらを統合するユニバースを構成していることを示した。すなわち、関係（スパン構造）を核として複数の意味の場がブリッジで緩やかに結ばれたこのユニバースの構造こそが、観測者の有無にかかわらず成り立つ「存在の普遍的な形式」にほかならない。私たちが経験する現象とは、この実在のユニバースの「内包」であり、観測から実在の全体構造を逆算することは原理的に不可能である。この知見は、ガブリエルが論じた「単一の包括的世界の不在」に数理的裏付けを与えると同時に、あらゆるモデル化技法において「単一の完璧な世界記述」を想定することの限界を示し、検証された前提のもとで複数のスキーマとブリッジのネットワークとして世界を捉える重要性を示唆している。

\begin{thebibliography}{99}

\bibitem{benabou1967}
B\'enabou, J. (1967). Introduction to Bicategories. In \textit{Lecture Notes in Mathematics} (Vol. 47, pp. 1--77). Springer.

\bibitem{fong2019}
Fong, B., \& Spivak, D. I. (2019). \textit{An Invitation to Applied Category Theory: Seven Sketches in Compositionality}. Cambridge University Press.（邦訳：『応用圏論』）

\bibitem{frege1892}
Frege, G. (1892). \"Uber Sinn und Bedeutung. \textit{Zeitschrift f\"ur Philosophie und philosophische Kritik}, 100, 25--50.（邦訳：「意義と意味について」）

\bibitem{gabriel2013}
Gabriel, M. (2013). \textit{Warum es die Welt nicht gibt}. Ullstein.（邦訳：『なぜ世界は存在しないのか』）

\bibitem{gabriel2015}
Gabriel, M. (2015). \textit{Fields of Sense: A New Realist Ontology}. Edinburgh University Press.

\bibitem{husserl1913}
Husserl, E. (1913). \textit{Ideen zu einer reinen Ph\"anomenologie und ph\"anomenologischen Philosophie. Erstes Buch: Allgemeine Einf\"uhrung in die reine Ph\"anomenologie}. Max Niemeyer.（邦訳：『イデーン I ― 純粋現象学と現象学的哲学のための諸構想 第1巻』）

\bibitem{kant1781}
Kant, I. (1781/1787). \textit{Kritik der reinen Vernunft}. Johann Friedrich Hartknoch.（邦訳：『純粋理性批判』）

\bibitem{lawvere1963}
Lawvere, F. W. (1963). Functorial Semantics of Algebraic Theories. \textit{Proceedings of the National Academy of Sciences}, 50(5), 869--872.

\bibitem{maclane1998}
Mac Lane, S. (1998). \textit{Categories for the Working Mathematician} (2nd ed.). Springer.（邦訳：『圏論の基礎』）

\bibitem{quine1969}
Quine, W. V. O. (1969). \textit{Ontological Relativity and Other Essays}. Columbia University Press.（邦訳：『存在論的相対性』）

\bibitem{spivak2012}
Spivak, D. I. (2012). Functorial Data Migration. \textit{Information and Computation}, 217, 31--51.

\bibitem{wittgenstein1921}
Wittgenstein, L. (1921). Logisch-Philosophische Abhandlung (Tractatus Logico-Philosophicus). \textit{Annalen der Naturphilosophie}.（邦訳：『論理哲学論考』）

\end{thebibliography}

\end{document}
